\chapter{Appendix}
\label{ch:appendix}

% Guideline 4.2.7: separate title page from the main body (the
% \chapter above + \appendix numbering in main.tex provides this),
% own numbering (A1, A1.1, ... below), source code referenced but not
% printed in full.

\section{QCar Hardware Safety Limits}
\label{app:safety-limits}

Table~\ref{tab:app-safety-limits} lists the documented hardware
safety limits enforced throughout every piece of drive-control code
presented in this thesis (Section~\ref{sec:qcar-platform}).

\begin{table}[htbp]
  \centering
  \caption{QCar documented hardware safety limits.}
  \label{tab:app-safety-limits}
  \begin{tabular}{@{}ll@{}}
    \toprule
    Limit & Value \\
    \midrule
    PWM duty cycle saturation & $\pm 30\%$ \\
    PWM duty cycle rate limit & $100\%/\mathrm{s}$ \\
    Battery auto-shutdown & below \SI{10.0}{\volt} \\
    Battery low-voltage warning & below \SI{10.5}{\volt} \\
    Steering command range & $\pm 0.5$~rad \\
    Max.\ physical steering angle & $\pm 0.5236$~rad ($\pm\ang{30}$) \\
    Overcurrent protection & \SI{5}{\ampere}/\SI{8}{\second}, \SI{10}{\ampere}/\SI{2}{\second}, or \SI{15}{\ampere}/\SI{0.5}{\second} $\to$ forced Neutral mode \\
    \bottomrule
  \end{tabular}
\end{table}

\section{Diagnostic and Calibration Tooling Reference}
\label{app:tooling}

Table~\ref{tab:app-tooling} summarises the diagnostic scripts
developed and used throughout Chapters~\ref{ch:hardware-bridge}
--\ref{ch:nav-integration}. Full source is maintained in the
project's version-controlled repository (electronic submission), not
reproduced here in full per the guideline's recommendation that
programming code be supplied on the electronic data carrier rather
than printed.

\begin{table}[htbp]
  \centering
  \caption{Diagnostic and calibration scripts referenced in this thesis.}
  \label{tab:app-tooling}
  \small
  \begin{tabular}{@{}p{3.8cm}p{9.5cm}@{}}
    \toprule
    Script & Purpose \\
    \midrule
    \texttt{calibrate\_steering.py} & Static steering-trim calibration aid (Section~\ref{sec:steering-trim}) \\
    \texttt{drive\_and\_log.py} & Drive to a target distance at a fixed speed, log \texttt{/odom} to CSV, exit automatically (Section~\ref{sec:final-calibration}) \\
    \texttt{check\_stop\_latency.py} & As above, plus reports cruise velocity and stop latency directly (Section~\ref{sec:deadband-discovery}) \\
    \texttt{monitor\_goal\_timing.py} & Instruments one Nav2 goal end-to-end: goal-sent, real/displayed motion start and stop, goal-reached (Section~\ref{sec:monitor-goal-timing}) \\
    \texttt{log\_goal.py} & Logs every \texttt{/plan} replan (sequence-numbered CSVs), \texttt{/odom}, and \texttt{/cmd\_vel} for one goal, for offline K-turn/cusp analysis (Section~\ref{sec:cusp-repro}) \\
    \texttt{\seqsplit{CuspStraightenerSmoother}} & \texttt{nav2\_core::Smoother} plugin resolving the reversal-cusp freeze (Section~\ref{sec:cusp-smoother-v2}); locked-in production configuration \\
    \texttt{\seqsplit{IsPathValidDebouncedCondition}} & Custom BT.CPP condition node requiring consecutive failed validity checks before replanning; implemented and hardware-tested, currently reverted/deselected (Section~\ref{sec:cusp-followon-reverted}) \\
    \texttt{\seqsplit{PathDeviationCritic}} & Custom MPPI critic pulling deviated trajectories toward a lookahead path point; implemented and tested, currently reverted/deselected (Section~\ref{sec:cusp-followon-reverted}) \\
    \bottomrule
  \end{tabular}
\end{table}

\section{Real-Hardware Bring-Up Command Reference}
\label{app:commands}

For reproducibility, the full bring-up sequence used throughout
Chapters~\ref{ch:hardware-bridge}--\ref{ch:nav-integration} is
reproduced here.

\begin{lstlisting}[language=bash,caption={Real-hardware bring-up (two QCar terminals, two development-PC terminals).}]
# --- On the QCar (two terminals) ---
sudo -E env PYTHONPATH=/home/nvidia/Documents/python \
    python3 ~/ros2_amit/src/onboard_bridge/qcar_bridge.py

PYTHONPATH=/home/nvidia/Documents/python \
    python3 ~/ros2_amit/src/onboard_bridge/qcar_lidar_node.py

# --- On the development PC ---
ros2 run qcar_updated qcar_relay_node.py --ros-args \
    -p qcar_ip:=<qcar-ip> -p qcar_clock_offset:=<measured-offset>

ros2 launch qcar_updated qcar_nav2.launch.py launch_sim:=false
\end{lstlisting}

\texttt{qcar\_clock\_offset} (Section~\ref{sec:clock-offset}) and the
QCar's DHCP-assigned IP address both drift session to session and
must be re-measured / re-checked before each bring-up, not reused
from a previous session.

\subsection{Safe Shutdown Sequence}
\label{app:shutdown}

The shutdown sequence established in Section~\ref{sec:safe-shutdown}
must be followed in order, using \texttt{SIGINT} (Ctrl-C) exclusively:

\begin{enumerate}
  \item Stop \texttt{qcar\_relay\_node.py} (development PC).
  \item Stop the visualisation launch process group (development PC).
  \item Stop \texttt{qcar\_lidar\_node.py} (QCar).
  \item Stop \texttt{qcar\_bridge.py} (QCar, requires \texttt{sudo}).
  \item Confirm both TCP ports are released on the QCar:
    \texttt{sudo ss -tlnp | grep -E "5555|5556"} should print nothing.
\end{enumerate}

\section{Source Code and Data Availability}
\label{app:source-availability}

% TODO: fill in with however your faculty/supervisor wants source
% code and raw data supplied - e.g. a git repository URL, or "supplied
% on the accompanying electronic data carrier as specified in
% guideline Section 4.2.7". If any data is restricted (see guideline
% Section 4.1.2), note that here instead.

The complete, version-controlled source code for the simulation
stack (Chapter~\ref{ch:sim-stack}), the hardware bridge
(Chapter~\ref{ch:hardware-bridge}), and all calibration/diagnostic
tooling (Appendix~\ref{app:tooling}) referenced throughout this
thesis is supplied on the accompanying electronic data carrier /
repository: \texttt{TODO: path or URL}.
